\documentclass[12pt]{article}
\usepackage[T1]{fontenc}
\usepackage[utf8]{inputenc}
\usepackage{lmodern}
\usepackage{textcomp}
\usepackage{enumitem}
\usepackage{geometry}
\geometry{margin=1in}
\begin{document}
\begin{center}
Answer Key - Sociology - Graduate Level Assessment 8 \
\small (Answers are ordered exactly as in the question paper.)
\end{center}

\section*{Multiple Choice}
\begin{enumerate}[label=\arabic*.]
\item \textbf{Answer:} D
\\ \textit{Options:} A) the world was divided into several rival overseas empires; B) industrialization provided better transport, technology and administration; C) there was intense economic, political and military competition; D) all of the above
\\ \textit{Note:} The correct answer "all of the above" is justified because various interrelated factors such as the rise of nationalism, the increase in industrial capabilities, and the formation of complex alliances collectively created a geopolitical environment that made war between nation states not only possible but also likely during the nineteenth century.
\item \textbf{Answer:} B
\\ \textit{Options:} A) the democratization of the family; B) putting an end to privatization; C) positive welfare: 'a hand up, not a hand-out'; D) the strengthening of civil society
\\ \textit{Note:} Giddens (1998) characterized New Labour's 'third way' as a blend of traditional social democratic principles with market-oriented reforms, emphasizing modernization and social inclusion rather than opposing privatization, which is why "putting an end to privatization" is not aligned with his definition.
\item \textbf{Answer:} D
\\ \textit{Options:} A) social solidarity and cohesion; B) the interpretive understanding of action; C) women's experiences and gendered knowledge; D) power, domination, and conflict
\\ \textit{Note:} Dahrendorf, Rex, and Habermas critically analyzed the structures of power and the dynamics of social conflict, emphasizing how domination and power relations shape social interactions and institutions.
\item \textbf{Answer:} D
\\ \textit{Options:} A) audiences selectively interpret what they want to hear; B) content analysis is the best way to identify the themes covered by the media; C) audiences passively absorb whatever messages they are given; D) social interaction reinforces the ideas and images that audiences select
\\ \textit{Note:} The media-themes model of media influence posits that through social interaction, individuals actively engage with and reinforce the ideologies and representations presented in media, thereby solidifying the ideas and images that they have selected for themselves.
\item \textbf{Answer:} B
\\ \textit{Options:} A) the division of each social class into the more or less privileged; B) a growing gap between rich and poor, resulting in class consciousness; C) the growth of intermediate strata in the middle classes; D) the tendency for the working class to live in very cold places
\\ \textit{Note:} The correct answer highlights Marx's theory that the increasing economic disparity between the bourgeoisie and the proletariat leads to heightened class consciousness, as the working class becomes aware of their exploitation and unified in their struggles against the wealthy elite.
\end{enumerate}

\section*{True / False}
\begin{enumerate}[label=\arabic*.]
\setcounter{enumi}{5}
\item \textbf{Answer:} True
\\ \textit{Options:} A) True; B) False
\\ \textit{Note:} Cultural restructuring involves revitalizing urban areas by integrating economic development strategies, reimagining industrial spaces for new purposes, and employing media advertising to reshape public perceptions and attract investment, thereby confirming the statement as true.
\item \textbf{Answer:} True
\\ \textit{Options:} A) True; B) False
\\ \textit{Note:} Bourdieu's theory posits that cultural capital, which includes knowledge, skills, and education, is crucial in maintaining and perpetuating social class distinctions by enabling individuals to navigate social structures and access opportunities that reinforce their social status.
\item \textbf{Answer:} False
\\ \textit{Options:} A) True; B) False
\\ \textit{Note:} Chodorow (1978) argued that girls develop stronger attachments to their mothers due to the mother-daughter relationship, while boys typically form a more complex attachment to their fathers as they navigate the process of separating from their mothers.
\item \textbf{Answer:} False
\\ \textit{Options:} A) True; B) False
\\ \textit{Note:} The 1944 Education Act established a framework for free and compulsory education for all children in England and Wales, ensuring that public schooling was accessible to lower-income families regardless of their financial situation.
\item \textbf{Answer:} False
\\ \textit{Options:} A) True; B) False
\\ \textit{Note:} The statement is false because both parents typically share responsibility for their child's welfare and socialization after a divorce, and external factors such as extended family, schools, and community resources also play significant roles in a child's development.
\end{enumerate}

\section*{Long Form Response}
\begin{enumerate}[label=\arabic*.]
\setcounter{enumi}{10}
\item \textbf{Answer:} The functionalist perspective views society as a complex system composed of interdependent parts, each serving a specific function that contributes to the overall stability and continuity of the social order. This framework, rooted in the works of sociologists like Émile Durkheim, posits that social institutions-such as family, education, and religion- work collaboratively to maintain societal equilibrium. Consequently, any disruption in one part of the system can lead to social instability, highlighting the importance of each component in sustaining societal health. However, functionalism also acknowledges that social change is inevitable and can arise from shifts in the functions of these parts, driven by factors such as technological advancements or shifts in cultural values. Understanding society through this lens enables sociologists to analyze both the mechanisms of stability and the conditions under which change occurs, thereby providing a comprehensive view of social dynamics.
\\ \textit{Note:} correct because it accurately describes the functionalist perspective's view of society as a system of interdependent parts that each fulfill specific functions, thereby emphasizing the importance of social institutions in maintaining stability and the potential consequences of disruptions within the system.
\item \textbf{Answer:} Murray's conception of the 'underclass' in sociology is characterized by a group of individuals who are marginalized, often residing in urban areas, and who experience chronic poverty, unemployment, and social isolation. He attributes to this group specific behavioral traits, including a reliance on welfare, single-parent households, and a perceived lack of motivation to improve their socioeconomic status. These characteristics have significant implications for social policy, as they suggest a need for interventions that address not only economic assistance but also family structure and cultural attitudes towards work. Moreover, Murray's theory raises critical questions within sociology regarding the interplay between individual agency and structural factors in perpetuating poverty. Thus, his conception of the underclass serves as both a descriptive framework and a call to reconsider the effectiveness of existing social policies.
\\ \textit{Note:} The answer correctly identifies Murray's characterization of the 'underclass' by emphasizing their socioeconomic challenges, behavioral traits, and the resultant implications for social policy, thereby demonstrating a comprehensive understanding of his sociological perspective.
\end{enumerate}

\vfill
\noindent\textit{End of Answer Key}
\end{document}